\section{Causal Logs and the Generation of Time}
\label{sec:time}

\subsection{History and Structure}

We distinguish two directed relations between nodes $i$ and $j$:

\begin{itemize}
  \item $H_{ij} \geq 0$: the \textbf{causal history} ---
    total accumulated influence from $i$ to $j$, irreversible and unfiltered.
  \item $C_{ij} \in [0,1)$: the \textbf{causal log} ---
    the selected portion of that history that becomes persistent structure.
    $C_{ij}$ represents the strength of retained influence:
    not what happened, but what continues to matter.
\end{itemize}

The fundamental distinction is:
\begin{equation}
  H_{ij} \neq C_{ij}.
  \label{eq:HneqC}
\end{equation}

The world is not the sum of everything that occurred,
but the structure of what was retained.

\subsection{The Self-Referential Selection Equation}

Selection is not externally imposed. It is generated by the structure itself.
Define the \textbf{log density} at node $i$:

\begin{equation}
  \rho_i = \sum_j C_{ij}.
  \label{eq:rho}
\end{equation}

The \textbf{causal log equation} is:

\begin{keybox}
\begin{equation}
  C_{ij} = 1 - \exp(-\lambda\,\rho_i\,H_{ij}),
  \label{eq:Cij}
\end{equation}
\end{keybox}

\noindent
where $\lambda > 0$ is the selection sensitivity, the sole external parameter.

Substituting \cref{eq:rho} into \cref{eq:Cij} yields a self-referential
fixed-point equation in $\rho_i$:

\[
  \rho_i = \sum_j \left[1 - \exp(-\lambda\,\rho_i\,H_{ij})\right].
\]

The stable fixed points $\rho_i^*$ of this equation correspond to
persistent physical structure.

\begin{remark}
The self-referential structure is not circular.
It is a fixed-point iteration: existing structure ($\rho_i$) governs
how new history ($H_{ij}$) is selected into new structure ($C_{ij}$).
Structure generates its own selectivity.
\end{remark}

\subsection{The Generation of Time}

Time is not a background parameter. It is defined as the rate at which
causal logs accumulate at a node:

\begin{keybox}
\begin{equation}
  \frac{d\tau_i}{dt} = \rho_i.
  \label{eq:tau}
\end{equation}
\end{keybox}

\noindent
Several consequences follow immediately:

\begin{itemize}
  \item Since $\rho_i \geq 0$ by construction, $\tau_i$ is monotonically
    non-decreasing. Irreversibility is built into the definition.
  \item Where $\rho_i = 0$, time does not advance.
  \item No separate thermodynamic argument for the arrow of time is required.
\end{itemize}

\begin{keybox}
\textit{Time is the rate at which structure accumulates.
Irreversibility is structural, not statistical.}
\end{keybox}

\subsection{The Log-Generation Rate}

Taking the time derivative of $\rho_i$ using \cref{eq:Cij} yields:

\begin{equation}
  \dot{\rho}_i = \frac{\lambda\,\rho_i\,B_i}{1 - \lambda\,A_i},
  \label{eq:rhodot}
\end{equation}

where $A_i = \sum_j e^{-\lambda\rho_i H_{ij}} H_{ij}$
and $B_i = \sum_j e^{-\lambda\rho_i H_{ij}} \dot{H}_{ij}$.
The numerator $\lambda\rho_i B_i$ measures the rate of new history injection;
the denominator $1 - \lambda A_i$ captures saturation by existing structure.
This microscopic expression underlies the log-generation rate $\sigma$
appearing in \cref{sec:flow}.
